% For next week:
% - Cover from section 7.4.1 to section 8.4 (included).
% - Complete the seventh iRAT by Thursday 27th at 13:00
% - Next class on Friday 28th at 11:00.


\subsection*{Question 1}
\textbf{Explain the two drawbacks of classical convergence analysis.}

The constants that determine the rate of convergence are expressed in terms of problem-dependent parameters (e.g., smoothness and strong convexity constants). These parameters are rarely known in practice.  

Additionally, the classical analysis relies on bounds that are coordinate-dependent. The problem parameters change with a change of coordinates, even though the actual convergence behavior of Newton's method remains unaffected. 

\subsection*{Question 2}
\textbf{Define self-concordance for a function $ f : \mathbb{R} \to \mathbb{R}$ and explain why it is coordinate invariant.}

A function \( f: \mathbb{R} \to \mathbb{R} \) is self-concordant if:  

\[
| f'''(x) | \leq 2 f''(x)^{3/2}, \quad x \in \textbf{dom }f
\]

It is coordinate invariant because it only involves second and third derivatives, which remains unchanged under affine transformations. 


\subsection*{Question 3}
\textbf{Consider a convex optimisation problem with equality constraints. Explain the method of reduction to unconstrained problems by elimination of the equality constraints.}

Given a convex optimization problem:

$$
\begin{array}{ll}
\min & f(x) \\
s.t. & Ax=b
\end{array}
$$

The feasible set can be parametrized as: 

$$
\{x:Ax=b\} = \{z\in\mathbb{R}^{n-p}:Fz + \hat x\}.
$$

Thus the problem is equivalent to: 

$$
\begin{array}{ll}
\displaystyle \min_z & f(Fz + \hat x).
\end{array}
$$

where the solution of the original problem is $x^* = Fz^* + \hat x$.

\subsection*{Question 4}
\textbf{Consider a convex optimisation problem with equality constraints. Explain the solution method based on the dual problem.}

The dual function of the convex optimisation problem is:

$$
g(\nu) = -b^\top \nu + \inf_x (f(x) + \nu^\top Ax) = -b^\top \nu - f^*(-A^\top \nu)
$$

where $f^*$ is the conjugate of $f$. So the dual problem of $(1)$ is

$$
\begin{array}{ll}
\max & -b^\top \nu - f^*(-A^\top \nu)
\end{array}
$$

Since by assumption there is an optimal point, the problem is strictly feasible, so Slater's condition holds. Therefore strong duality holds, and the dual optimum is attained. 

If $g$ is twice differentiable, the optimal solution $\nu^*$ can be found solving $(5)$ using the unconstrained minimisation method. The optimal primal solution $x^*$ can then be found by solving $\nabla L(x, \nu^*) = 0$. 


\subsection*{Question 5}
\textbf{Explain why the infeasible start Newton's method is not a descent method.}


The infeasible satrt Newton's method is not a descent method as it does not always guarantee that the function value decreases monotonically at each step. The derivative:

$$
\left.\frac{d}{dt} f(x+t\Delta x)\right|_{t=0} = \nabla f(x)^\top \Delta x = - \Delta x^\top (\Delta^2 f(x)\Delta x + A^\top w) = - \Delta x^\top \Delta^2 f(x)\Delta x + (Ax-b)^\top w
$$

is not necessarily negative unless $x$ is feasible and so $Ax-b=0$.



